\lecture

\begin{guiding}
Linear algebra has functors: duals, homomorphisms, tensor products. Do they act
on vector bundles? And can we attach to every bundle algebraic invariants ---
cohomology classes --- that detect twisting?
\end{guiding}

\section{Continuous functors of vector bundles}

\begin{definition}
Let $\Theta$ be the category whose objects are finite dimensional real vector
spaces and whose morphisms are isomorphisms. A functor
\[
T\colon \Theta\times\dots\times\Theta\longrightarrow\Theta
\]
of $k$ variables is \emph{continuous} if $T(f_1,\dots,f_k)$ depends continuously
on $f_1,\dots,f_k$, the latter being regarded as matrices.
\end{definition}

\begin{theorem}\label{thm:functor-bundle}
Let $T$ be a continuous functor of $k$ variables and let $\xi_1,\dots,\xi_k$ be
vector bundles over $B$. Then there is a vector bundle
$T(\xi_1,\dots,\xi_k)$ over $B$ whose fibre over $b$ is
$T\big(F_b(\xi_1),\dots,F_b(\xi_k)\big)$.
\end{theorem}

\begin{proof}
\begin{strategy}
The set and the fibres are prescribed; what has to be produced is a topology.
Local trivializations of the $\xi_i$ give candidate charts, and the only thing
to check is that the change of chart is continuous --- which is exactly the
continuity hypothesis on $T$.
\end{strategy}

\step{1} \emph{The set.} Put
\[
E=\bigsqcup_{b\in B}T\big(F_b(\xi_1),\dots,F_b(\xi_k)\big),
\]
with the obvious projection $\pi$ to $B$ and the vector space structure of
$T(\dots)$ on each fibre.

\step{2} \emph{Candidate charts.} Suppose $U\subset B$ trivializes all the
$\xi_i$ simultaneously, with trivializations giving, for $b\in U$, isomorphisms
$h_{ib}\colon\R^{n_i}\to F_b(\xi_i)$ depending continuously on $b$. Since $T$
is a functor, applying it to these isomorphisms gives a linear isomorphism
\[
T(h_{1b},\dots,h_{kb})\colon
T(\R^{n_1},\dots,\R^{n_k})\longrightarrow
T\big(F_b(\xi_1),\dots,F_b(\xi_k)\big),
\]
and hence a bijection
\[
H_U\colon U\times T(\R^{n_1},\dots,\R^{n_k})\longrightarrow\pi^{-1}(U),
\qquad
(b,x)\longmapsto T(h_{1b},\dots,h_{kb})(x).
\]

\step{3} \emph{Compatibility of charts.} Let $U_1,U_2$ be two such open sets.
On $U_1\cap U_2$ the composite $H_{U_1}^{-1}\circ H_{U_2}$ is
\[
(b,x)\longmapsto\Big(b,\ T\big(h^{(1)}_{1b}{}^{-1}h^{(2)}_{1b},\ \dots,\
h^{(1)}_{kb}{}^{-1}h^{(2)}_{kb}\big)(x)\Big),
\]
using functoriality. Each $h^{(1)}_{ib}{}^{-1}h^{(2)}_{ib}$ is the transition
isomorphism of $\xi_i$, which depends continuously on $b$; by continuity of $T$
the whole expression depends continuously on $(b,x)$. The same applies to the
inverse.

\step{4} \emph{The topology.} Declare a subset of $E$ open when its preimage
under every $H_U$ is open. By Step 3 the charts $H_U$ are then homeomorphisms
onto their images, and they are fibrewise linear isomorphisms by construction.
Hence $E$ is a vector bundle with the prescribed fibres.
\end{proof}

\begin{example}
The functor $V\mapsto\Hom(V,\R)$ yields the \emph{dual bundle} $\xi^*$; two
variables give $\xi\otimes\eta$, $\Hom(\xi,\eta)$, $\xi\oplus\eta$, and so on.
\end{example}

\begin{remark}
If $\xi$ is Euclidean, the metric gives an isomorphism $\xi\iso\xi^*$, namely
$(b,v)\mapsto(b,\langle v,\cdot\rangle_b)$; it is continuous and fibrewise an
isomorphism, so Lemma~\ref{lem:fibrewise-iso} applies. Over a paracompact base a
Euclidean metric always exists, by patching local metrics with a partition of
unity.
\end{remark}

\section{The Stiefel--Whitney axioms}

We now \emph{postulate} a theory of invariants and derive consequences; the
construction of the classes will occupy Lectures 4--6. This shortcut lets us
prove interesting theorems immediately, which is also how the subject developed
historically.

\begin{definition}[Axioms]\label{def:sw-axioms}
To each $\R^n$-bundle $\xi$ over a base $B$ there are associated cohomology
classes
\[
w_i(\xi)\in H^i(B;\Z_2),\qquad i\ge0,
\]
the \emph{Stiefel--Whitney classes}, subject to:

\begin{enumerate}
\item[(A1)] \textbf{Rank.} $w_0(\xi)=1$ and $w_i(\xi)=0$ for $i>n=\rk\xi$.

\item[(A2)] \textbf{Naturality.} If $\hat f\colon\xi\to\eta$ is a bundle map
with base map $f$, then $w_i(\xi)=f^*w_i(\eta)$; equivalently
$w_i(f^*\eta)=f^*w_i(\eta)$.

\item[(A3)] \textbf{Whitney product formula.} If $\xi,\eta$ have the same base,
\[
w_k(\xi\oplus\eta)=\sum_{i=0}^{k}w_i(\xi)\smile w_{k-i}(\eta).
\]

\item[(A4)] \textbf{Normalization.} $w_1(\gamma^1_1)\neq0$ for the canonical
line bundle $\gamma^1_1\to\RP^1$.
\end{enumerate}
\end{definition}

\begin{proposition}
If $\xi\iso\eta$ then $w_i(\xi)=w_i(\eta)$ for all $i$.
\end{proposition}

\begin{proof}
An isomorphism is a bundle map whose base map is the identity, so (A2) gives
$w_i(\xi)=\mathrm{id}^*w_i(\eta)=w_i(\eta)$.
\end{proof}

\begin{proposition}\label{prop:trivial-vanishes}
If $\xi$ is trivial then $w_i(\xi)=0$ for every $i>0$.
\end{proposition}

\begin{proof}
\begin{strategy}
A trivial bundle is pulled back from a bundle over a point, and a point has no
cohomology in positive degrees.
\end{strategy}

\step{1} Write $\xi\iso B\times\R^n$ and let $\eps'$ denote the bundle
$\mathrm{pt}\times\R^n$ over a point. The map
$B\times\R^n\to\mathrm{pt}\times\R^n$, $(b,x)\mapsto(\mathrm{pt},x)$, is a
bundle map whose base map is the constant map $c\colon B\to\mathrm{pt}$.

\step{2} By (A2), $w_i(\xi)=c^*w_i(\eps')$. Since
$H^i(\mathrm{pt};\Z_2)=0$ for $i>0$, we have $w_i(\eps')=0$ for $i>0$, hence
$w_i(\xi)=0$.
\end{proof}

\begin{proposition}[Stability]\label{prop:stability}
If $\eps$ is a trivial bundle then $w_i(\eps\oplus\eta)=w_i(\eta)$.
\end{proposition}

\begin{proof}
By Proposition~\ref{prop:trivial-vanishes}, $w_0(\eps)=1$ and $w_j(\eps)=0$ for
$j>0$; the Whitney formula (A3) then reduces to the single term
$w_0(\eps)\smile w_i(\eta)=w_i(\eta)$.
\end{proof}

\begin{proposition}\label{prop:sections-kill-top}
If a Euclidean $\R^n$-bundle $\xi$ admits $k$ nowhere dependent sections, then
\[
w_{n}(\xi)=w_{n-1}(\xi)=\dots=w_{n-k+1}(\xi)=0 .
\]
\end{proposition}

\begin{proof}
\begin{strategy}
The sections span a trivial sub-bundle; splitting it off drops the rank by $k$,
and the rank axiom does the rest.
\end{strategy}

\step{1} Let $s_1,\dots,s_k$ be the sections and let $\eps^k\subset\xi$ be the
sub-bundle whose fibre over $b$ is $\spann(s_1(b),\dots,s_k(b))$. It is a
sub-bundle of rank $k$, trivialized by
$(b,x)\mapsto\sum x_i s_i(b)$ as in Theorem~\ref{thm:trivial-iff-sections}.

\step{2} Let $\eta=(\eps^k)^\perp$, a bundle of rank $n-k$ with
$\eps^k\oplus\eta\iso\xi$ by Theorem~\ref{thm:orthogonal-complement}.

\step{3} By Proposition~\ref{prop:stability},
$w_i(\xi)=w_i(\eps^k\oplus\eta)=w_i(\eta)$, and by (A1) applied to $\eta$ this
vanishes for $i>n-k$.
\end{proof}

\begin{remark}
Read contrapositively: if $w_{n-k}(TM)\neq0$ then $M$ admits at most $k$
everywhere independent vector fields. Stiefel--Whitney classes are
\emph{obstructions to finding independent sections}, and this is the historical
origin of the theory.
\end{remark}

\section{The total class and Whitney duality}

Let $H^{\Pi}(B;\Z_2)$ denote the ring of formal series
$a=a_0+a_1+a_2+\cdots$ with $a_i\in H^i(B;\Z_2)$, with the evident sum and
product; the product is commutative and associative. The \emph{total
Stiefel--Whitney class} of an $\R^n$-bundle is
\[
w(\xi)=1+w_1(\xi)+\dots+w_n(\xi)\ \in\ H^{\Pi}(B;\Z_2),
\]
and the Whitney formula (A3) becomes the single identity
\[
w(\xi\oplus\eta)=w(\xi)\,w(\eta).
\]

\begin{proposition}\label{prop:invertible}
The series in $H^\Pi(B;\Z_2)$ with leading term $1$ form an abelian group under
multiplication.
\end{proposition}

\begin{proof}
\begin{strategy}
Solve $w\bar w=1$ degree by degree; each equation determines the next unknown
uniquely.
\end{strategy}

\step{1} Given $w=1+w_1+w_2+\cdots$, we seek $\bar w=1+\bar w_1+\bar w_2+\cdots$
with $w\bar w=1$. Comparing terms of degree $k\ge1$ gives
\[
\bar w_k+\bar w_{k-1}w_1+\dots+\bar w_1 w_{k-1}+w_k=0 .
\]

\step{2} This determines $\bar w_k$ uniquely in terms of $\bar w_1,\dots,
\bar w_{k-1}$ and $w_1,\dots,w_k$, so by induction $\bar w$ exists and is
unique. The first values are
\[
\bar w_1=w_1,
\qquad
\bar w_2=w_1^2+w_2,
\qquad
\bar w_3=w_1^3+w_3 .
\]

\step{3} Associativity and commutativity are inherited from
$H^*(B;\Z_2)$, and the identity element is $1$.
\end{proof}

\begin{lemma}[Whitney duality]\label{lem:whitney-duality}
If $M\subset\R^N$ is a smooth submanifold with normal bundle $\nu$, then
\[
w(\nu)=\bar w(TM),
\qquad\text{that is,}\qquad
w(TM)\,w(\nu)=1 .
\]
\end{lemma}

\begin{proof}
By Example~\ref{ex:normal-bundle}, $TM\oplus\nu\iso T\R^N|_M$, which is trivial.
Proposition~\ref{prop:trivial-vanishes} gives $w(TM\oplus\nu)=1$, and the
Whitney formula turns this into $w(TM)w(\nu)=1$; now apply
Proposition~\ref{prop:invertible}.
\end{proof}

\begin{example}
The normal bundle of $S^n\subset\R^{n+1}$ is trivial, spanned by the outward
unit normal, so $w(\nu)=1$ and therefore $w(TS^n)=1$: all Stiefel--Whitney
classes of spheres vanish. Yet $TS^2$ is non-trivial. The classes are powerful,
but they are not complete invariants.
\end{example}

\section{Computations on projective spaces}

Recall $H^*(\RP^n;\Z_2)=\Z_2[a]/(a^{n+1})$ with $a$ the generator of
$H^1(\RP^n;\Z_2)$.

\begin{example}\label{ex:w-gamma}
$w(\gaman)=1+a$.

Indeed, the inclusion $j\colon\RP^1\hookrightarrow\RP^n$ induced by
$\R^2\subset\R^{n+1}$ is covered by a bundle map
$\gamma^1_1\to\gaman$, so (A2) gives $w_1(\gamma^1_1)=j^*w_1(\gaman)$. By (A4)
the left side is non-zero, so $w_1(\gaman)\neq0$; since
$H^1(\RP^n;\Z_2)=\Z_2\{a\}$, this forces $w_1(\gaman)=a$. The bundle has rank
$1$, so $w_i=0$ for $i\ge2$ by (A1).
\end{example}

\begin{example}\label{ex:w-gamma-perp}
Let $\gamma^\perp$ be the orthogonal complement of $\gaman$ inside the trivial
bundle $\RP^n\times\R^{n+1}$. Then
\[
w(\gamma^\perp)=1+a+a^2+\dots+a^n .
\]
Indeed $\gaman\oplus\gamma^\perp$ is trivial, so
$w(\gamma^\perp)=\overline{(1+a)}$, and one checks directly in
$\Z_2[a]/(a^{n+1})$ that
\[
(1+a)(1+a+\dots+a^n)=1+a^{n+1}=1 .
\]
\end{example}

\begin{lemma}\label{lem:TRPn}
There are isomorphisms
\[
T\RP^n\ \iso\ \Hom(\gaman,\gamma^\perp)
\qquad\text{and}\qquad
T\RP^n\oplus\eps^1\ \iso\
\underbrace{\gaman\oplus\dots\oplus\gaman}_{n+1},
\]
where $\eps^1$ is the trivial line bundle.
\end{lemma}

\begin{proof}
\begin{strategy}
Describe tangent vectors of $\RP^n$ as antipodal pairs of tangent vectors of
$S^n$, and observe that such a pair is exactly a linear map from the line
through $x$ to its orthogonal complement. The second isomorphism then comes from
adding the summand $\Hom(\gaman,\gaman)$, which is trivial, and using
$\gamma^\perp\oplus\gaman=\eps^{n+1}$.
\end{strategy}

\step{1} \emph{Tangent vectors upstairs and downstairs.} Let $q\colon S^n\to
\RP^n$ be the quotient map. Since $q$ is a local diffeomorphism, $Dq$ identifies
\[
T\RP^n=\big\{\,\{(x,v),(-x,-v)\}\ \big|\ x\cdot x=1,\ x\cdot v=0\,\big\},
\]
because $Dq(x,v)=Dq(-x,-v)$ and these are the only identifications.

\step{2} \emph{From pairs to homomorphisms.} Given such a pair, let
$L=\R x\subset\R^{n+1}$ be the corresponding line, that is, the fibre of $\gaman$
over $\{\pm x\}$, and define a linear map
\[
\ell\colon L\longrightarrow L^\perp,
\qquad
\ell(x)=v .
\]
Replacing $(x,v)$ by $(-x,-v)$ gives the same $\ell$, since
$\ell(-x)=-v$. Conversely $\ell$ determines the pair. This assignment is a
bijection
\[
T\RP^n\longrightarrow \Hom(\gaman,\gamma^\perp),
\]
continuous with continuous inverse in local trivializations, and linear on each
fibre; Lemma~\ref{lem:fibrewise-iso} makes it an isomorphism.

\step{3} \emph{The trivial summand.} The line bundle
$\Hom(\gaman,\gaman)$ has the nowhere zero section $b\mapsto\mathrm{id}$, hence
is trivial by Theorem~\ref{thm:trivial-iff-sections}:
$\Hom(\gaman,\gaman)\iso\eps^1$.

\step{4} \emph{Adding up.} Using Step 2, Step 3 and additivity of $\Hom$ in the
second variable,
\[
T\RP^n\oplus\eps^1
\iso\Hom(\gaman,\gamma^\perp)\oplus\Hom(\gaman,\gaman)
\iso\Hom\big(\gaman,\ \gamma^\perp\oplus\gaman\big),
\]
and the right-hand side is $\Hom\big(\gaman,\eps^{n+1}\big)$,
since $\gamma^\perp\oplus\gaman$ is the trivial bundle
$\RP^n\times\R^{n+1}$.

\step{5} \emph{Splitting the trivial bundle.} Finally
\[
\Hom\big(\gaman,\eps^{n+1}\big)
\iso\bigoplus_{i=1}^{n+1}\Hom\big(\gaman,\eps^1\big)
=\bigoplus_{i=1}^{n+1}(\gaman)^*
\iso\bigoplus_{i=1}^{n+1}\gaman,
\]
the last step because a line bundle with a Euclidean metric is isomorphic to its
own dual.
\end{proof}

\begin{theorem}\label{thm:w-RPn}
$w(T\RP^n)=(1+a)^{n+1}$, that is,
$w_i(T\RP^n)=\binom{n+1}{i}a^i\bmod 2$.
\end{theorem}

\begin{proof}
By Proposition~\ref{prop:stability} and Lemma~\ref{lem:TRPn},
\[
w(T\RP^n)=w\big(T\RP^n\oplus\eps^1\big)
=w\big((\gaman)^{\oplus(n+1)}\big)
=w(\gaman)^{n+1}=(1+a)^{n+1},
\]
using the Whitney formula and Example~\ref{ex:w-gamma}.
\end{proof}

The mod-2 Pascal triangle
\[
\begin{array}{ccccccccc}
 & & & &1& & & &\\
 & & &1& &1& & &\\
 & &1& &0& &1& &\\
 &1& &1& &1& &1&\\
1& &0& &0& &0& &1
\end{array}
\]
already gives $w(T\RP^3)=1$ and $w(T\RP^4)=1+a+a^4$.

\begin{corollary}[Stiefel]\label{cor:parallelizable}
If $\RP^n$ is parallelizable then $n+1$ is a power of $2$.
\end{corollary}

\begin{proof}
\begin{strategy}
Parallelizable means $T\RP^n$ trivial, hence $w=1$. Over $\Z_2$ the Frobenius
identity $(1+a)^{2^r}=1+a^{2^r}$ lets us factor $(1+a)^{n+1}$ and exhibit a
surviving term unless $n+1$ is a power of $2$.
\end{strategy}

\step{1} If $T\RP^n$ is trivial then $w(T\RP^n)=1$ by
Proposition~\ref{prop:trivial-vanishes}, so $(1+a)^{n+1}=1$ in
$\Z_2[a]/(a^{n+1})$.

\step{2} Write $n+1=2^r m$ with $m$ odd. Since squaring is additive mod $2$,
\[
(1+a)^{2^r}=1+a^{2^r},
\]
and therefore
\[
(1+a)^{n+1}=\big(1+a^{2^r}\big)^{m}
=1+m\,a^{2^r}+\binom{m}{2}a^{2\cdot 2^r}+\cdots
\]

\step{3} Suppose $m>1$. Then $2^r<n+1$, so $2^r\le n$ and the class $a^{2^r}$ is
non-zero in $\Z_2[a]/(a^{n+1})$. Its coefficient is $m\equiv1\bmod2$, so
$(1+a)^{n+1}\neq1$, contradicting Step 1. Hence $m=1$ and $n+1=2^r$.
\end{proof}

\begin{remark}
The converse is false: the only parallelizable real projective spaces are
$\RP^1,\RP^3,\RP^7$, corresponding to $\C$, the quaternions and the octonions.
This requires deeper tools (Bott--Milnor, Kervaire, Adams) and is not proved
here.
\end{remark}

\begin{example}[Immersions of $\RP^9$]\label{ex:immersion-RP9}
Suppose $M^n$ immerses in $\R^{n+k}$. By Lemma~\ref{lem:normal-immersion} the
normal bundle $\nu_f$ has rank $k$, and $TM\oplus\nu_f$ is trivial, so that
$w_i(\nu_f)=\bar w_i(TM)$. Therefore
\[
\bar w_i(TM)=0\qquad\text{for } i>k .
\]
For $M=\RP^9$: since $10=8+2$,
\[
w(T\RP^9)=(1+a)^{10}=(1+a^{8})(1+a^{2})=1+a^{2}+a^{8}
\]
in $\Z_2[a]/(a^{10})$. Inverting, with $b=a^2$ and $b^5=a^{10}=0$, the equation
$(1+b+b^4)\,\bar w=1$ gives successively the coefficients $1,1,1,1,0$, so
\[
\bar w(T\RP^9)=1+a^{2}+a^{4}+a^{6}.
\]
The top non-zero term sits in degree $6$, so $k\ge6$: the space $\RP^9$ cannot be
immersed in $\R^{9+k}$ for $k<6$.
\end{example}

\begin{example}[Immersions of $\RP^{2^r}$]
Let $n=2^r$. Then
\[
w(T\RP^n)=(1+a)^{2^r+1}=\big(1+a^{2^r}\big)(1+a)=1+a+a^{n}
\]
in $\Z_2[a]/(a^{n+1})$, and one checks
\[
\bar w(T\RP^{n})=1+a+a^2+\dots+a^{n-1},
\]
since $(1+a)(1+a+\dots+a^{n-1})=1+a^n$ telescopes and
$a^{n}(1+a+\dots+a^{n-1})=a^n$ modulo $a^{n+1}$, so the two extra terms cancel.
The top non-zero term is $a^{n-1}$, so any immersion of $\RP^{2^r}$ needs
$k\ge n-1$, that is, an ambient $\R^{2n-1}$. This is exactly the bound of
Whitney's immersion theorem, so projective spaces of $2$-power dimension are the
worst case.
\end{example}

\section{Stiefel--Whitney numbers}

\begin{guiding}
Classes live in the cohomology of $M$, which makes them awkward to compare
between different manifolds. Can we extract from them plain \emph{numbers}, and
what do those numbers see?
\end{guiding}

\begin{definition}
Let $M$ be a closed smooth $n$-manifold with fundamental class
$\mu_M\in H_n(M;\Z_2)$. For every partition
$r_1+2r_2+\dots+nr_n=n$ the associated \emph{Stiefel--Whitney number} is
\[
w_1^{r_1}\cdots w_n^{r_n}[M]
:=\big\langle\, w_1(TM)^{r_1}\smile\cdots\smile w_n(TM)^{r_n},\ \mu_M\,\big\rangle
\ \in\ \Z_2 .
\]
Two $n$-manifolds \emph{have the same Stiefel--Whitney numbers} if these agree
for every such $(r_1,\dots,r_n)$.
\end{definition}

\begin{example}
For $n$ even, $w_n(T\RP^n)=\binom{n+1}{n}a^n=(n+1)a^n\neq0$ and likewise
$w_1(T\RP^n)^n=(n+1)^na^n\neq0$, so $\RP^n$ has non-zero Stiefel--Whitney
numbers. For $n$ odd they all vanish, as we verify at the start of the next
lecture --- and this is no accident: odd projective spaces are boundaries.
\end{example}

In Lecture 3 we prove that Stiefel--Whitney numbers vanish on boundaries, and
state Thom's theorem that they classify closed manifolds up to unoriented
cobordism.
